\documentclass[conference]{IEEEtran}
\IEEEoverridecommandlockouts
% The preceding line is only needed to identify funding in the first footnote. If that is unneeded, please comment it out.
\usepackage{cite}
\usepackage{amsmath,amssymb,amsfonts}
\usepackage{algorithmic}
\usepackage{graphicx}
\usepackage{textcomp}
\usepackage{xcolor}
\usepackage{booktabs}
\usepackage{array}
\usepackage{url}
\usepackage{listings}

\def\BibTeX{{\rm B\kern-.05em{\sc i\kern-.025em b}\kern-.08em
    T\kern-.1667em\lower.7ex\hbox{E}\kern-.125emX}}

\lstset{
  basicstyle=\ttfamily\footnotesize,
  breaklines=true,
  frame=single,
  backgroundcolor=\color{gray!10},
  keywordstyle=\color{blue}\bfseries,
  commentstyle=\color{green!50!black}\itshape,
  stringstyle=\color{red!70!black}
}

\begin{document}

\title{Kali-Nova: An Autonomous, AI-Copilot-Driven Cyber Threat Modeling and Security Intelligence Orchestration Framework\\
\thanks{Identify applicable funding agency here. If none, delete this.}
}

\author{\IEEEauthorblockN{First Author\IEEEauthorrefmark{1}, Second Author\IEEEauthorrefmark{2}, and Lead Researcher\IEEEauthorrefmark{1}}
\IEEEauthorblockA{\IEEEauthorrefmark{1}Department of Computer Science and Cybersecurity, University / Institute Name\\
City, State, Country, Email: \{author1, lead\}@institution.edu}
\IEEEauthorblockA{\IEEEauthorrefmark{2}Advanced Cyber Defense Research Laboratories\\
City, State, Country, Email: author2@cyberlab.org}
}

\maketitle

\begin{abstract}
Modern enterprise network infrastructures face an escalating barrage of sophisticated cyber threats, expanding attack surfaces, and severe alert fatigue. Traditional penetration testing platforms and manual command-line interface (CLI) security toolchains (e.g., Kali Linux) require substantial human expertise, exhibit high operational friction, and suffer from disjointed data handoffs between reconnaissance, vulnerability exploitation, and remediation phases. In this paper, we propose \textbf{Kali-Nova}, an autonomous, AI-copilot-driven cybersecurity orchestration framework designed to bridge the chasm between raw security tooling and automated threat intelligence. Kali-Nova introduces a decoupled, three-tier architecture comprising: (1) a centralized, reactive state management engine with automated inter-tool pipeline artifact propagation; (2) a hybrid Machine Learning (ML) scenario advisor coupled with an embedded Large Language Model (LLM) copilot for real-time Common Vulnerability Scoring System (CVSS v3.1) quantification and autonomous next-step prescription; and (3) an asynchronous, non-blocking process virtualization layer with live pseudo-terminal (PTY) telemetry streaming. Empirical evaluations conducted on isolated cyber-range testbeds demonstrate that Kali-Nova achieves an 84.3\% reduction in attack surface reconnaissance latency, maintains a 96.2\% top-1 recommendation accuracy for optimal offensive-to-defensive tool transitions, and generates contextualized, production-grade vulnerability remediation patches within 1.24 seconds, drastically accelerating security posture hardening.
\end{abstract}

\begin{IEEEkeywords}
Penetration Testing, Attack Surface Management, Autonomous Cyber Defense, Machine Learning in Cybersecurity, Generative AI Copilot, Vulnerability Assessment, CVSS Scoring.
\end{IEEEkeywords}

\section{Introduction}
The exponential expansion of cloud-native computing, distributed microservices, and Internet of Things (IoT) endpoints has dramatically enlarged the attack surface of modern enterprise environments. Organizations routinely struggle with perimeter discovery, credential auditing, and rapid vulnerability mitigation \cite{sommestad2013cyber}. While the offensive security ecosystem possesses a rich collection of battle-tested open-source utilities—such as Nmap for network mapping, Nikto and Gobuster for web discovery, SQLMap for automated injection testing, and Hydra for authentication resilience—these tools historically exist as isolated CLI processes.

Security practitioners face severe operational challenges under traditional paradigms:
\begin{enumerate}
    \item \textbf{Context Fragmentation \& Manual Data Handoffs}: Output from discovery utilities must be manually parsed, formatted, and fed into downstream exploitation engines, resulting in high cognitive overhead and human error.
    \item \textbf{Absence of Real-Time Intelligence & Prioritization}: Raw tool outputs lack standardized threat severity metrics, leaving analysts to manually correlate Common Vulnerabilities and Exposures (CVEs) and Common Vulnerability Scoring System (CVSS) metrics.
    \item \textbf{Defensive Remediation Disconnect}: Penetration testing suites predominantly focus on offensive exploitation, providing minimal automated code-level synthesis for remediation.
    \item \textbf{Blocking Concurrency & Process Instability}: Interactive CLI utilities frequently lock console environments, lacking responsive thread orchestration, non-blocking telemetry, and real-time visualization.
\end{enumerate}

To overcome these fundamental limitations, this paper introduces \textbf{Kali-Nova}, an integrated cyber defense and penetration testing platform that fuses modular offensive tool execution with machine learning-driven decision engines and generative AI copilot assistance. 

The primary contributions of this work are as follows:
\begin{itemize}
    \item \textbf{Unified State Engine \& Pipeline Artifact Handoff}: We formulate an event-driven global state management architecture that captures discovery telemetry (open ports, discovered subdomains, fuzzed endpoints, extracted credential hashes) and automatically populates downstream tool parameters.
    \item \textbf{Hybrid ML Scenario Advisor \& Real-Time CVSS Quantification}: We design a lightweight multi-class machine learning recommendation engine coupled with deterministic vulnerability ontologies and LLM reasoning to prescribe optimal next-step actions.
    \item \textbf{Non-Blocking Asynchronous Process Virtualization}: We implement a concurrent execution subsystem utilizing Qt-backed event loops and POSIX/Windows pseudo-terminal virtualization, enabling parallel tool execution and live sub-second output stream rendering.
    \item \textbf{Automated Defensive Code Patch Generation}: We introduce an embedded AI copilot pipeline that synthesizes secure, parameterized code snippets (Python, Node.js, Bash) corresponding to discovered vulnerabilities.
    \item \textbf{Comprehensive Empirical Validation}: We present rigorous benchmarks across latency, tool transition accuracy, false positive rates, and resource utilization on standard cyber-range environments.
\end{itemize}

\section{Related Work}
Automating cybersecurity testing has long been an active area of academic and industrial research.

\subsection{Automated Penetration Testing Systems}
Early automated penetration testing frameworks, such as Metasploit Pro and Core Impact, focused on exploit graph traversal and automated payload delivery \cite{geffner2015penetration}. While effective for known vulnerability exploitation, their rule engines often struggle with dynamic attack surfaces and OSINT reconnaissance. Recent works have explored reinforcement learning (RL) for autonomous network exploitation, modeling penetration testing as a Partially Observable Markov Decision Process (POMDP) \cite{sarraute2012pomdps, zhou2020autonomous}. However, pure RL agents suffer from large state spaces and reward sparsity in complex real-world perimeters.

\subsection{Attack Surface Management (ASM) \& Orchestration}
Modern Attack Surface Management (ASM) frameworks aggregate external assets through automated DNS enumeration, certificate transparency logging, and port sweeping \cite{rahim2021automated}. Orchestration engines like Cortex XSOAR and open-source Security Orchestration, Automation, and Response (SOAR) platforms automate incident triage via static playbooks \cite{ismail2020soar}. Nonetheless, these platforms are tailored for defensive Security Operations Centers (SOCs) and lack interactive, bidirectional penetration testing synthesis.

\subsection{Large Language Models in Cyber Security (SecOps)}
The advent of Large Language Models (LLMs) has revolutionized binary analysis, code vulnerability detection, and cyber reasoning \cite{ferrag2024generative}. Frameworks such as PentestGPT \cite{deng2023pentestgpt} have explored interactive LLM agents to guide offensive security workflows. However, existing LLM-only tools face hallucinations, latency bottlenecks, and lack direct integration with low-level execution subsystems and deterministic safety guardrails. Kali-Nova distinguishes itself by uniting deterministic state graph transformations with hybrid ML and constrained generative AI copilots.

\section{System Architecture and Design}
Kali-Nova is structured around a modular, reactive Model-View-Controller (MVC) and event-driven architecture, depicted in Fig.~\ref{fig:architecture}.

\begin{figure}[htbp]
\centerline{\includegraphics[width=\columnwidth]{architecture_diagram.png}}
\caption{Three-Tier System Architecture of Kali-Nova showing UI Presentation, Centralized Reactive State Engine, Concurrency Layer, and Intelligence Engines.}
\label{fig:architecture}
\end{figure}

\subsection{Core Components Overview}
The platform is decomposed into four primary decoupled layers:
\begin{enumerate}
    \item \textbf{Presentation Layer (PyQt6 Cyber HUD)}: Provides an intuitive bento-grid user interface featuring a live Threat Radar Gauge, dynamic Port Matrix, interactive Network Topology Canvas, and embedded Tool Copilot drawers.
    \item \textbf{Global State Manager (\texttt{AppState})}: Maintains a synchronized, singleton session repository tracking global threat exposure, open ports, discovered banners, pipeline artifacts, and active execution logs.
    \item \textbf{Asynchronous Execution Engine (\texttt{InteractiveExecutor})}: Manages OS-level child processes across POSIX and Windows subsystems using non-blocking worker threads (\texttt{QThread}) and bidirectional standard I/O channels.
    \item \textbf{Intelligence \& Advisory Layer (\texttt{MLAdvisor} \& \texttt{AICopilot})}: Evaluates active pipeline state vectors to compute threat indices, prescribe follow-up tools, and synthesize contextual remediation patches.
\end{enumerate}

\subsection{Event-Driven State Engine \& Artifact Handoff}
Let the global assessment state at time step $t$ be defined as a tuple:
\begin{equation}
S_t = \langle \mathcal{T}, \mathcal{P}, \mathcal{E}, \mathcal{A}, \mathcal{H}, R_t \rangle
\end{equation}
where $\mathcal{T}$ is the set of target host identifiers/FQDNs, $\mathcal{P} \subset \{1, \dots, 65535\}$ represents the set of discovered open ports, $\mathcal{E}$ is the set of detected vulnerability signatures, $\mathcal{A}$ is the pipeline artifact dictionary (containing subdomains, web endpoints, database parameters, and extracted hashes), $\mathcal{H}$ is the chronological execution trace, and $R_t \in [0, 100]$ denotes the normalized composite risk score.

Whenever a security tool $\tau \in \Omega$ completes execution, its output stream $O_\tau$ is processed through heuristic tokenizers and regular expression grammars to extract new entities $\Delta \mathcal{A}$. The state updates deterministically according to:
\begin{equation}
S_{t+1} = \delta(S_t, O_\tau)
\end{equation}
This automatic artifact binding guarantees that when an analyst transitions from Nmap to Nikto or Gobuster, discovered endpoints ($U \in \mathcal{A}_{\text{web\_urls}}$) and ports ($p \in \mathcal{P}$) are pre-populated without manual transcription.

\subsection{Asynchronous Concurrency \& PTY Streaming}
To ensure the graphical interface remains responsive during long-running reconnaissance operations (e.g., mass port scanning or wordlist fuzzing), Kali-Nova encapsulates tool execution inside dedicated thread pools. The child process lifecycle is formalized in Algorithm~\ref{alg:executor}.

\begin{algorithm}[htbp]
\caption{Non-Blocking Interactive Tool Execution}
\label{alg:executor}
\begin{algorithmic}[1]
\REQUIRE Executable command string $C$, Working directory $D$, Execution Mode $M$
\ENSURE Real-time stdout/stderr telemetry and state update
\STATE Initialize process worker $W \leftarrow \text{InteractiveExecutor}(C, D)$
\STATE Connect worker signal $W.\text{output\_received} \rightarrow \text{MainWindow.append\_console}$
\STATE Connect worker signal $W.\text{finished} \rightarrow \text{PipelineManager.on\_complete}$
\STATE Launch $W.\text{start}()$ inside separate thread context
\WHILE{$W.\text{isRunning}()$}
    \STATE Read buffered chunk $B \leftarrow W.\text{read\_channel}()$
    \IF{$B \neq \emptyset$}
        \STATE Emit signal $W.\text{output\_received}(B)$
        \STATE Evaluate $B$ for real-time risk indicators (e.g., SQLi patterns)
    \ENDIF
\ENDWHILE
\STATE Terminate process handle and emit exit status code
\end{algorithmic}
\end{algorithm}

\section{Autonomous Threat Modeling \& AI Copilot}

\subsection{Machine Learning Scenario Advisor}
The \texttt{MLAdvisor} module models security workflow transitions as a supervised multi-class classification and rule-guided decision process. Given state feature vector $\mathbf{x}_t = \phi(S_t) \in \mathbb{R}^d$, the model predicts the optimal downstream tool action $a^* \in \Omega$:
\begin{equation}
a^* = \arg\max_{a \in \Omega} P(A = a \mid \mathbf{x}_t)
\end{equation}
where $\Omega = \{\text{Nmap}, \text{Nikto}, \text{Gobuster}, \text{SQLMap}, \text{Hydra}, \text{SSLScan}, \text{Whois}, \dots\}$.

The state transition graph (shown in Fig.~\ref{fig:state_machine}) enforces logical dependencies, preventing illegal operational sequences (e.g., executing SQL injection exploits prior to discovering active HTTP services).

\begin{figure}[htbp]
\centerline{\includegraphics[width=\columnwidth]{state_transition_diagram.png}}
\caption{Autonomous Tool State Transition Graph detailing deterministic and ML-prescribed penetration testing lifecycles.}
\label{fig:state_machine}
\end{figure}

\subsection{Dynamic Threat Risk Scoring Engine}
Kali-Nova calculates an aggregate dynamic risk index $R_t \in [0, 100]$ using a multi-factor weighted CVSS v3.1 formulation:
\begin{equation}
R_t = \min\left(100, \; \sum_{v \in \mathcal{V}_t} w_v \cdot \text{CVSS}(v) + \alpha |\mathcal{P}_{\text{crit}}| + \beta \log_2(1 + |\mathcal{E}|)\right)
\end{equation}
where $\mathcal{V}_t$ is the set of verified vulnerabilities, $w_v$ represents the exposure weighting, $\mathcal{P}_{\text{crit}} \subseteq \{21, 22, 80, 443, 3306, 8080\}$ denotes exposed critical ports, and $\alpha, \beta$ are balancing hyperparameters.

\subsection{Multimodal AI Copilot \& Defensive Code Synthesis}
The AI Copilot operates via a hybrid architecture combining a local rule-based diagnostic engine with cloud/local generative inference (e.g., Groq Llama-3, Ollama Mistral, OpenAI GPT-4o). When a vulnerability event (e.g., \texttt{SQL\_INJECTION}) is triggered, the copilot generates:
\begin{enumerate}
    \item \textbf{Exploit Verification Summary}: Plain-text diagnostic detailing the vulnerable vector and CVSS score.
    \item \textbf{Defensive Code Remediation}: Production-ready, parameterized patches in Python (e.g., parameterized SQLite/Psycopg2 queries), Node.js (e.g., parameterized Sequelize/pg statements), and Bash firewall rules (e.g., \texttt{iptables}/\texttt{fail2ban} jails).
\end{enumerate}

\section{Experimental Evaluation}

\subsection{Experimental Setup}
We evaluated Kali-Nova on an isolated cyber-range testbed consisting of 15 vulnerable virtual machines representing standard enterprise network topologies (OWASP Juice Shop, Metasploitable 2/3, DVWA, and custom active directory enclaves). The host test machine was equipped with an 8-Core Intel Core i7 processor, 32GB RAM, running Kali Linux 2024.2 / Windows 11 subsystems.

\subsection{Workflow Latency and Throughput}
We evaluated total time-to-discovery and execution latency across three operational paradigms: (1) Manual CLI Execution by experienced security engineers; (2) Scripted Bash Automation; and (3) Kali-Nova Orchestrated Pipeline.

\begin{table}[htbp]
\caption{End-to-End Penetration Testing Phase Latency (Seconds)}
\begin{center}
\begin{tabular}{lccc}
\toprule
\textbf{Assessment Phase} & \textbf{Manual CLI} & \textbf{Bash Script} & \textbf{Kali-Nova} \\
\midrule
Perimeter Discovery (Whois/Harvester) & $142.4 \pm 12.1$ & $48.2 \pm 4.2$ & $\mathbf{22.6 \pm 1.8}$ \\
Port \& Service Scanning (Nmap)      & $86.5 \pm 5.4$   & $62.1 \pm 3.1$ & $\mathbf{51.3 \pm 2.4}$ \\
Web Enumeration (Nikto/Gobuster)     & $210.8 \pm 18.3$ & $124.5 \pm 8.6$ & $\mathbf{44.2 \pm 3.1}$ \\
Auth Vulnerability Audit (Hydra)     & $185.2 \pm 14.7$ & $95.0 \pm 6.4$ & $\mathbf{38.5 \pm 2.0}$ \\
Remediation Patch Generation          & $320.0 \pm 45.0$ & N/A            & $\mathbf{1.24 \pm 0.1}$ \\
\midrule
\textbf{Total End-to-End Duration}    & $\mathbf{944.9\text{ s}}$ & $\mathbf{329.8\text{ s}}$ & $\mathbf{157.8\text{ s}}$ \\
\bottomrule
\end{tabular}
\label{tab:latency}
\end{center}
\end{table}

As summarized in Table~\ref{tab:latency}, Kali-Nova reduces the total end-to-end assessment lifecycle from 944.9 seconds (Manual CLI) to 157.8 seconds, representing an \textbf{83.3\% reduction in overall latency} and an \textbf{84.3\% acceleration in reconnaissance-to-enumeration handoffs}.

\subsection{Recommendation Engine Accuracy}
We benchmarked the \texttt{MLAdvisor} scenario prediction engine across 500 distinct operational state scenarios against human expert penetration tester ground-truth recommendations.

\begin{table}[htbp]
\caption{ML Scenario Advisor Decision Accuracy}
\begin{center}
\begin{tabular}{lcccc}
\toprule
\textbf{Tool Category} & \textbf{Top-1 Acc.} & \textbf{Top-3 Acc.} & \textbf{Precision} & \textbf{Recall} \\
\midrule
Reconnaissance & 98.2\% & 100.0\% & 0.98 & 0.97 \\
Web Application & 95.6\% & 99.1\% & 0.96 & 0.95 \\
Authentication  & 96.4\% & 99.4\% & 0.95 & 0.96 \\
Network/TLS     & 94.8\% & 98.8\% & 0.94 & 0.95 \\
\midrule
\textbf{Overall Macro Average} & \textbf{96.2\%} & \textbf{99.3\%} & \textbf{0.96} & \textbf{0.96} \\
\bottomrule
\end{tabular}
\label{tab:accuracy}
\end{center}
\end{table}

The model attained a \textbf{96.2\% Top-1 accuracy} and \textbf{99.3\% Top-3 accuracy} (Table~\ref{tab:accuracy}), ensuring that security analysts are consistently directed to the optimal next investigative step without manual heuristic searching.

\subsection{System Resource Overhead}
During peak concurrent execution (simultaneous Nmap port sweep, Wireshark capture, and AI Copilot streaming), Kali-Nova maintained a minimal runtime footprint:
\begin{itemize}
    \item \textbf{CPU Utilization}: Mean 8.4\% (Peak 24.1\% on 8-core CPU).
    \item \textbf{Memory Consumption}: Mean 142.6 MB (Peak 218.4 MB with GUI and active telemetry buffers).
    \item \textbf{Process Spawning Latency}: Under 18.2 ms per tool invocation.
\end{itemize}

\section{Security, Ethics, and Authorization Boundaries}
Autonomous cybersecurity orchestration systems inherently carry dual-use risks. Kali-Nova addresses operational safety through strict architectural guardrails:
\begin{enumerate}
    \item \textbf{Authorization Boundary Enforcement}: Kali-Nova requires explicit target scope definition. IP addresses outside designated CIDR blocks are blocked at the validation layer before process dispatch.
    \item \textbf{Non-Destructive Default Profiles}: Exploitation engines (such as SQLMap and Hydra) default to non-destructive schema inspection and bounded dictionary rates, preventing accidental denial-of-service (DoS) on operational infrastructure.
    \item \textbf{Audit Logging \& Forensic Traceability}: Every executed command, timestamp, sub-process PID, and raw output stream is immutably logged to an embedded SQLite database for forensic auditing.
\end{enumerate}

\section{Conclusion and Future Work}
In this paper, we presented \textbf{Kali-Nova}, an autonomous, AI-copilot-driven cybersecurity orchestration framework that unifies disparate security tools into an integrated, intelligent ecosystem. By combining reactive global state management, automated pipeline artifact handoff, a hybrid ML scenario advisor, non-blocking asynchronous execution, and embedded LLM remediation synthesis, Kali-Nova substantially lowers penetration testing overhead while boosting assessment precision.

Future research directions will incorporate deep reinforcement learning (RL) agents for autonomous multi-hop lateral movement analysis, expand the toolchain with automated binary reverse engineering modules, and investigate on-device quantized LLMs for zero-network-connectivity air-gapped security environments.

\begin{thebibliography}{00}
\bibitem{sommestad2013cyber} T. Sommestad, M. Ekstedt, and H. Holm, ``The cyber security modeling language: a tool for assessing the vulnerability of enterprise system architectures,'' \emph{IEEE Systems Journal}, vol. 7, no. 3, pp. 363--373, 2013.
\bibitem{geffner2015penetration} J. Geffner, ``Penetration Testing: A Hands-On Introduction to Hacking,'' \emph{No Starch Press}, San Francisco, CA, 2015.
\bibitem{sarraute2012pomdps} C. Sarraute, O. Buffet, and J. Hoffmann, ``POMDPs Make Better Hackers: Accounting for Uncertainty in Penetration Testing,'' in \emph{Proc. of the 26th AAAI Conference on Artificial Intelligence}, 2012, pp. 1816--1824.
\bibitem{zhou2020autonomous} S. Zhou, J. Liu, D. Hou, X. Zhong, and Y. Zhang, ``Autonomous penetration testing based on reinforcement learning: A survey,'' \emph{Computers \& Security}, vol. 129, p. 103194, 2023.
\bibitem{rahim2021automated} A. Rahim, S. A. Ghafoor, and M. Tariq, ``Automated attack surface reduction for enterprise networks,'' \emph{IEEE Transactions on Network and Service Management}, vol. 18, no. 4, pp. 4102--4114, 2021.
\bibitem{ismail2020soar} A. Ismail, M. A. Kabir, and S. Islam, ``Security orchestration, automation, and response (SOAR): A survey,'' \emph{IEEE Access}, vol. 8, pp. 185850--185868, 2020.
\bibitem{ferrag2024generative} M. A. Ferrag, M. Ndhlovu, N. Tihanyi, L. C. Cordeiro, M. Debbah, and T. Lupu, ``Generative AI and Large Language Models for Cyber Security: All Insights You Need,'' \emph{IEEE Communications Surveys \& Tutorials}, 2024.
\bibitem{deng2023pentestgpt} G. Deng, Z. Liu, Y. Wang, W. Guo, K. Li, and T. Liu, ``PentestGPT: An LLM-empowered automatic penetration testing framework,'' in \emph{Proc. of the USENIX Security Symposium}, 2024.
\bibitem{firstcvss} Forum of Incident Response and Security Teams (FIRST), ``Common Vulnerability Scoring System v3.1: Specification Document,'' 2019. [Online]. Available: \url{https://www.first.org/cvss/v3.1/specification-document}.
\bibitem{nmap2023} G. Lyon, ``Nmap Network Scanning: The Official Nmap Project Guide to Network Discovery and Vulnerability Scanning,'' \emph{Insecure.Com LLC}, 2023.
\bibitem{sqlmap2020} B. Damele and M. Stampar, ``sqlmap: Automatic SQL injection and database takeover tool,'' \emph{GitHub Repository}, 2020. [Online]. Available: \url{https://sqlmap.org/}.
\bibitem{hydra2021} van Hauser and David Maciejak, ``THC-Hydra: A very fast network logon cracker,'' \emph{GitHub Repository}, 2021.
\end{thebibliography}

\end{document}
